\subsection{Dinâmica orbital}
\label{sec:dinamica_orbital}
Segundo \cite{wie2008space}, a dinâmica orbital é análise do movimento orbital de espaçonaves sob influencia gravitacional e outras forças externas.
As Leis de Newton do movimento e a Lei da gravidade, assim como as três Leis do movimento planetário de Kepler providenciam o fundamento para esse estudo.

% \subsubsection{Problema dos dois corpos}
Nesta dissertação foi considerado o problema mais simples da mecânica orbital, ou seja, o problema dinâmico de duas massas pontuais sob a influência de sua atração gravitacional mútua, conhecido como o problema dos dois corpos na mecânica celeste, sendo este utilizado para descrever o movimento absoluto das espaçonaves em missão de \gls{rvdb}.
A formulação do problema matemático é descrita pelas leis de movimento e pela lei da gravidade de Newton \cite{wie2008space}. O problema de dois corpos será formulado com as três leis empíricas do movimento planetário de Kepler, que são:

\begin{enumerate}[label=(\roman*)]
    \item Cada planeta move-se ao redor do Sol descrevendo uma órbita elíptica, tendo o Sol localizado em um dos focos da elipse (Lei das Órbitas).
    \item O vetor raio traçado do Sol até o planeta varre áreas iguais em intervalos de tempo iguais, conservando a quantidade de movimento angular.
    \item O quadrado do período orbital de um planeta é proporcional ao cubo do semi-eixo maior de sua órbita, relação que expressa a dependência entre o tempo de revolução e a distância média ao Sol (Lei dos Períodos).
\end{enumerate}

Conforme \cite{wie2008space}, este problema de dois corpos é ainda mais simplificado para o movimento quando se considera a dinâmica orbital de espaçonaves porque a massa do corpo primário (planeta) é muito maior que a do corpo secundário (espaçonave), ou seja, $m_1 \gg m_2$, o que resulta em $\mu \approx Gm_1$.
Portanto, para um sistema Terra--espaçonave, temos $\mu \approx \mu_\oplus = G M_\oplus$, onde $\mu_\oplus$ denota o parâmetro gravitacional da Terra e $G M_\oplus$ denota a massa da Terra.
Vale ressaltar que o corpo primário não é fixado inercialmente no problema dos dois corpos.

O problema dos dois corpos deve ser diferenciado do chamado problema restrito de dois corpos (sistema planeta--satélite), no qual o corpo primário de massa $m_1$ é assumido fixo inercialmente. Esse problema restrito de dois corpos é frequentemente descrito pelo movimento de força central de uma partícula de massa $m_2$ em torno de um corpo primário de massa $m_1$ fixado inercialmente.
Nesse estudo, então, o sistema de dois corpos é o sistema restrito dos dois corpos \cite{wie2008space}.

\subsubsection{O problema dos corpos: Movimento absoluto e relativo}

% Como exemplo prático do problema dos dois corpos, pode-se citar a órbita da Lua em torno da Terra, bem como a de um planeta em torno do Sol \cite{wgsantos2011}.
A aceleração de cada corpo é dada por:

\begin{equation}
\label{eq:wie3_1}
m_1 \ddot{\vec R}_1 =
+ G \frac{m_1m_2}{r^3} \vect r
\end{equation}

\begin{equation}
\label{eq:wie3_2}
m_2 \ddot{\vec R}_2 =
- G \frac{m_1m_2}{r^3} \vect r
\end{equation}

onde:
\begin{equation}
    \label{eq:vetor_posicao_r}
\vect{r} = \vec R_{\, 2} - \vec R_{\, 1}.
\end{equation}

A equação \eqref{eq:vetor_posicao_r} é o vetor posição de $m_2$ em relação a $m_1$ e em relação ao centro da Terra \cite{wie2008space}; e, $r$ é a distância entre um corpo em relação ao outro. Portanto, somando as equações \eqref{eq:wie3_1} e \eqref{eq:wie3_2}, obtém-se:

\begin{equation}
\label{eq:somaAceleracao_Corpo1e2}
m_1 \ddot{\vec R}_1 + m_2 \ddot{\vec R}_2 = 0
\end{equation}

Definindo o vetor posição $\vec R_{c}$ para o centro de massa do sistema como

\begin{equation}
\vec R_{c} =
\frac{m_1\,\vec {R_1} + m_2\,\vec R_2}{m_1+m_2},
\end{equation}

portanto,

\begin{equation}
 \ddot{\vec R}_{c} = 0.
\end{equation}

Integrando a equação anterior, obtém-se

\begin{equation}
\dot{\vec{R}}_c(t) = \dot{\vec{R}}_c(0) = \text{vetor constante},
\end{equation}

\begin{equation}
\vec{R}_c(t) = \vec{R}_c(0) + \dot{\vec{R}}_c(0)\,t,
\end{equation}

onde $\vec{R}_{c}(0)$ e $\dot{\vec{R}}_{c}(0)$ representam a posição e a velocidade inercial do centro de massa no instante $t = 0$.
Assim, metade dos 12 estados do sistema dinâmico foi resolvida. Os seis estados restantes podem ser obtidos considerando o movimento relativo entre as duas massas.
Subtraindo a \eqref{eq:wie3_1} multiplicada por $m_2$ da \eqref{eq:wie3_2} multiplicada por $m_1$, obtém-se

\begin{equation}
\ddot{\vec{r}} + \mu\,\frac{\vec{r}}{r^3} = 0,
\end{equation}

sendo reescrito como
\begin{equation}
\label{eq:eqFundamentalDoisCorpos}
\ddot {\vec r} = -\frac{\mu}{\vec {r^3}}\vec r
\end{equation}

onde $\ddot{\vect r}$ é a aceleração inercial de $m_2$ em relação a $m_1$; $\mu$ é o parâmetro gravitacional do sistema de dois corpos, e \eqref{eq:eqFundamentalDoisCorpos} é a equação fundamental (geral) do problema dos dois corpos \cite{fehse2003,wie2008space}.
A dinâmica orbital do alvo (sem aceleração de controle) é propagada por\fn{O subescrito \emph{p} se refere à espaçonave perseguidora, enquanto que \emph{a} é a espaçonave alvo.}

\begin{equation}
\label{eq:movimentoalvo}
\ddot{\vec r}_a = \frac{-\mu}{\vec {r^3_a}}\vec {r}_a.
\end{equation}

Enquanto que o perseguidor (com aceleração de controle) é propagado por

\begin{equation}
\label{eq:movimentoperseguidor}
\ddot{\vec r}_p
=
-\frac{\mu}{\vec {r_{p}^3}}\vec r_{p}
+
\frac{\vec F}{m_{p}},
\end{equation}

onde $\vect F$ é um tipo de força não-gravitacional, como uma força de controle, aplicada no perseguidor para se aproximar do alvo.
A Figura~\eqref{fig:movimento_ref_absoluto} apresenta o modelo físico da \gls{emtr} no contexto do movimento descrito no referencial inercial geocêntrico.
Nesse sistema, as posições do alvo e da espaçonave perseguidora são expressas em relação ao centro da Terra, possibilitando a análise do movimento absoluto.

\begin{figure}[ht]
\centering
\includegraphics[width=0.9\linewidth]{3 Formulacao matematica/Imagens/Posicao absoluta.png}
\caption{O problema do movimento relativo formulado para missões de \gls{rvdb}.}
\label{fig:movimento_ref_absoluto}
\end{figure}

% \newpage